\section*{Chapter \ref{ch:qm:ito-calculus} --- Itô Calculus}

\begin{solution}{exo:qm:ito-calculus:1}
$\E\int_0^1W_t^2\,dt = \int_0^1t\,dt = \tfrac12$. With \cref{ex:qm:ito-calculus:wdw},
$\Var(\tfrac12(W_1^2 - 1)) = \tfrac14\Var(W_1^2) = \tfrac14 \times 2 = \tfrac12$, and the mean is zero.
\end{solution}

\begin{solution}{exo:qm:ito-calculus:2}
$d(W^3) = 3W^2\,dW + 3W\,dt$. Taking expectations, $\E[W_t^3] = 3\int_0^t\E[W_s]\,ds = 0$, and
$W_t^3 - 3\int_0^tW_s\,ds$ is a \omterm{def:qm:probability-at-speed:martingale}{martingale}.
\end{solution}

\begin{solution}{exo:qm:ito-calculus:3}
Mean $e^{1} = 2.72$, median $e^{(0.10 - 0.08) \times 10} = 1.22$. The median is flat when $\sigma^2 = 2\mu$,
$\sigma = \sqrt{0.2} = 44.7\%$.
\end{solution}

\begin{solution}{exo:qm:ito-calculus:4}
It is $\mathcal E(W)_t$: $dZ = Z\,dW$, and $\E\int_0^tZ_s^2\,ds = \int_0^te^s\,ds < \infty$, so the \omterm{def:qm:ito-calculus:integral}{Itô integral}
is a true \omterm{def:qm:probability-at-speed:martingale}{martingale}. $\E[e^{2W_1}] = e^{2}= 7.39$ (the moment generating function of $\mathcal
N(0,1)$ at 2).
\end{solution}

\begin{solution}{exo:qm:ito-calculus:5}
$d(tW_t) = W_t\,dt + t\,dW_t$, so $\int_0^1W_t\,dt = W_1 - \int_0^1t\,dW_t = \int_0^1(1 - t)\,dW_t$, with variance
$\int_0^1(1 - t)^2dt = \tfrac13$ by the isometry.
\end{solution}

\begin{solution}{exo:qm:ito-calculus:6}
$[X, Y]_t = 2t$, $[Y]_t = 5t$, correlation $2/\sqrt5 = 0.894$.
\end{solution}

\begin{solution}{exo:qm:ito-calculus:7}
$L = 60$ gives $a = 0.223$ and a look-ahead Sharpe ratio of 3.45; the average over 100 seeded
walks is 3.47. A slower average puts less weight on the last increment, so the bias is smaller,
but it is still enormous.
\end{solution}

\begin{solution}{exo:qm:ito-calculus:8}
The signal uses the closing price, which is not known until the closing auction in which the
rule claims to trade has printed: the decision is not adapted to the information available
when the order must be sent. Compute the signal from the auction's indicative price published
before the order deadline (chapter~1), or trade at the next open.
\end{solution}

\begin{solution}{pb:qm:ito-calculus:1}
\textbf{1.} $c = 0.01\,(\sum_{m=1}^{19}(m/20)^2)^{1/2} = 0.01 \times 2.485 = 0.0248$.
\textbf{2.} $a = 0.95/2.485 = 0.382$.
\textbf{3.} Mean $a\sigma = 0.0038$ and standard deviation $\sigma\sqrt{1 + a^2} = 0.0107$ per day, per
unit of position.
\textbf{4.} $\sqrt{252} \times 0.382/1.070 = 5.67$.
\textbf{5.} Zero: $\theta_{t}$ is independent of $S_{t+1} - S_t$.
\textbf{6.} 5.43 same-bar, $-0.04$ next-bar.
\textbf{7.} Same-bar $5.69 \pm 0.13$; next-bar $0.00 \pm 0.32$.
\textbf{8.} Both equal 9.653 (965\% of notional); the difference is $2 \times 10^{-14}$.
\textbf{9.} 34.9.
\textbf{10.} With no edge the annualised Sharpe ratio estimated over $T$ years has a standard error
of about $1/\sqrt T$ (chapter~11): $0.32$ for ten years.
\textbf{11.} $\sum\theta_t\Delta S_t = \sum\theta_{t-1}\Delta S_t + \sum\Delta\theta_t\Delta S_t$: an Itô sum plus the
covariation $[\theta, S]$.
\textbf{12.} Its integrand is adapted, so it is a \omterm{def:qm:probability-at-speed:martingale}{martingale} started at zero.
\textbf{13.} 9.36, 7.47 and 3.45.
\textbf{14.} The daily Sharpe ratio per bar is unchanged, and annualising over $252 \times 78$ bars
multiplies it by $\sqrt{78}$: 50.
\textbf{15.} Yes, if the signal uses only the close and the gains start at the next open; the
position then misses the overnight move, which the honest backtest must also leave out.
\textbf{16.} Shift the positions by one bar and see whether the performance collapses; run
\texttt{lookahead\_test}, or regress the daily P\&L on $\Delta\theta_t\Delta S_t$; and read the
timestamps of the signal's inputs.
\textbf{17.} No: its sign is the sign of $[\theta, S]$; a mean-reversion rule, which sells as the price
rises, has a negative look-ahead bias.
\textbf{18.} The bounce between bid and ask makes last-trade prices mean-revert; an honest
mean-reversion rule on those prices ``earns'' the bounce, which costs the spread to capture
(One Quant Book~10).
\textbf{19.} Named result: \emph{the look-ahead in the backtest}: on a random walk the same-bar
backtest of a unit-variance rule with correlation $a$ to the latest increment has a daily Sharpe
ratio $a/\sqrt{1 + a^2}$; for the 20-day trend rule, $a = 0.382$ and the annualised Sharpe ratio is 5.67
(5.69 in simulation), against zero for the Itô sum.
\textbf{20.} Each position is fixed from information available before the price move it is
multiplied by.
\end{solution}

\begin{solution}{iq:qm:ito-calculus:1}
$\tfrac12(W_T^2 - T)$. The left-point sums differ from $\tfrac12\sum\Delta(W^2)$ by $\tfrac12\sum(\Delta W)^2$,
which tends to $T$: the \omterm{def:qm:brownian-motion:qv}{quadratic variation} does not vanish.
\iqlookfor{the telescoping identity and the \omterm{def:qm:brownian-motion:qv}{quadratic variation}.}
\end{solution}

\begin{solution}{iq:qm:ito-calculus:2}
$\mu - \sigma^2/2 = 2\%$ a year. The mean of $S_{10}/S_0$ is $e^{1} = 2.72$ but the median is $1.22$, and a
tenth of holders end below 0.24: the mean is carried by a few large outcomes.
\iqlookfor{Itô's correction and mean versus median.}
\end{solution}

\begin{solution}{iq:qm:ito-calculus:3}
No: $d(W^2) = 2W\,dW + dt$, so $W_t^2 - t$ is the \omterm{def:qm:probability-at-speed:martingale}{martingale}. $\Var(\int_0^TW\,dW) = \int_0^Tt\,dt =
T^2/2$.
\iqlookfor{Itô's formula and the isometry.}
\end{solution}

\begin{solution}{iq:qm:ito-calculus:4}
Drive it by events in time order: every datum carries the time it became available, the
strategy sees an as-of view, and an order it emits at $t$ can fill only against prices after
$t$ plus a latency. The engine, not the strategy, applies positions to the next interval, and
a test compares every run with a shifted run and flags covariation between positions and
concurrent returns.
\iqlookfor{availability timestamps, as-of joins, and a structural shift.}
\end{solution}

\begin{solution}{iq:qm:ito-calculus:5}
A process that is a \omterm{def:qm:probability-at-speed:martingale}{martingale} only up to a sequence of \omterm{def:qm:probability-at-speed:stopping}{stopping times}; a positive one is a
\omterm{def:qm:probability-at-speed:martingale}{supermartingale} and may lose expectation, like $1/|B_t|$ for a three-dimensional \omterm{def:qm:brownian-motion:bm}{Brownian motion}.
A candidate density for a \omterm{def:qm:probability-at-speed:change}{change of measure}, or a deflated price, that is a strict \omterm{def:qm:ito-calculus:local}{local
martingale} gives prices that violate put--call parity or a measure that is not a probability.
\iqlookfor{localisation, Fatou, and the pricing consequence.}
\end{solution}

\begin{solution}{iq:qm:ito-calculus:6}
Finance uses Itô: a position must be set before the price move, which is the left point, and
the resulting gains are \omterm{def:qm:probability-at-speed:martingale}{martingales} under the right measure. Physics often uses Stratonovich,
the limit of smooth noise, because it keeps the ordinary chain rule. They differ by $\tfrac12[X,W]$.
\iqlookfor{non-anticipation versus smooth-noise limits.}
\end{solution}
